В диссертационной работе рассмотрены методы безопасного обучения с подкреплением для задач, где агенту необходимо не только достигать цели, но и ограничивать штрафы за небезопасные действия во время обучения и исполнения. Такой конфликт особенно заметен в робототехнических задачах: для нахождения успешного поведения агент должен исследовать среду, однако именно исследование часто приводит к столкновениям, попаданию в опасные области или другим нарушениям ограничений.

В работе последовательно рассмотрены три аспекта этой проблемы. Первый связан с использованием моделей мира для предварительной оценки последствий действий. Второй~--- с иерархической декомпозицией поведения, которая позволяет работать с задачами большого горизонта через промежуточные подцели. Третий~--- с переходом к постановке, где цель и ограничение безопасности задаются не только численно, но и на естественном языке.

Проведённый обзор показывает, что существующие методы решают эти задачи только частично. Одноуровневые безопасные стратегии могут снижать штрафы, но в задачах большого горизонта часто становятся слишком осторожными и перестают продвигаться к цели. Методы на основе моделей мира дают возможность оценивать риск до выполнения действия, однако их качество зависит от точности модели, горизонта предсказания и выбранного способа оценки штрафов. Подходы с априорной информацией о среде могут давать более жёсткий контроль безопасности, но обычно требуют известной динамики, явных функций ограничений или доступа к информации, которая недоступна стандартному агенту обучения с подкреплением.

В рамках направления безопасного обучения с подкреплением на основе модели мира исследован подход SPOWL~\citeinsynopsis{latyshev2025spowl}, в котором модель мира используется для оценки награды, штрафов и долгосрочного риска кандидатных действий. Показано, что безопасное исследование должно учитывать не только мгновенный штраф, но и оценки будущих накопленных штрафов, позволяющие оценивать риск за пределами короткого горизонта планирования. Экспериментальный анализ показал, что способ агрегации ансамбля функций оценки будущих накопленных штрафов существенно влияет на баланс между безопасностью и эффективностью: минимальная агрегация недооценивает риск, максимальная приводит к избыточной консервативности, а усреднение обеспечивает наиболее устойчивое поведение агента.

Для решения проблемы безопасного исследования в задачах с большим горизонтом разработан метод SG-Safe~\citeinsynopsis{gorbov2026sgsafe}. Метод использует стратегию генерации промежуточных подцелей, которая применяется на этапе обучения и направляет безопасную стратегию к более информативным областям пространства состояний. При этом на этапе исполнения используется только итоговая безопасная стратегия, поэтому метод не требует онлайн-планирования и дополнительных вычислительных затрат. Эксперименты на Safety Gym и POLAMP показали, что SG-Safe повышает успешность решения навигационных задач при сохранении низкого уровня нарушений безопасности. В среднем на Safety Gym метод достигает доли успешных эпизодов $0.87$ при частоте столкновений $0.03$, а на POLAMP демонстрирует преимущество в задачах с длинными маршрутами и узкими коридорами.

Дальнейшее развитие иерархического подхода реализовано в методе ITES~\citeinsynopsis{gorbov2026ites}. В отличие от SG-Safe, метод ITES вводит явный контроль безопасности как на уровне генерации подцелей, так и на уровне их достижения. Верхнеуровневая стратегия генерирует промежуточные подцели с учётом модели штрафов и высокоуровневого критика штрафов, а низкоуровневая стратегия использует модель мира для проверки безопасности коротких воображаемых траекторий. Эксперименты на SafeAntMaze и SafePusher показали, что ITES превосходит неиерархические методы безопасного обучения с подкреплением в задачах с большим горизонтом, где многие методы сравнения либо нарушают ограничения, либо не достигают цели. В SafeAntMaze C-shape метод достигает доли успешных эпизодов $0.88$ при штрафов $15.2$, а в SafePusher~--- доли успешных эпизодов $0.42$ при штрафов $12.8$, сохраняя ограничения безопасности.

Отдельно в диссертации рассмотрена задача безопасного обучения с подкреплением с естественно-языковыми целями и ограничениями. Для этой постановки разработан бенчмарк Caged CrafText~\citeinsynopsis{gorbov2026caged}, расширяющий CrafText до задач, где агент получает не только инструкцию цели, но и текстовое ограничение безопасности. Бенчмарк включает бюджетные, последовательные, математические и отношенческие ограничения, проверяющие память, безопасное исследование, рассуждение, связывание инструкции с объектами и состояниями среды, а также долгосрочное планирование. Эксперименты с PPO, PPO-LAG, PPO-Cost и агентом на основе большой языковой модели показали, что существующие подходы не решают задачу устойчивого совмещения следования инструкциям и безопасности. PPO достигает награды, но систематически нарушает ограничения; PPO-LAG снижает штрафы, но часто становится чрезмерно консервативным; агент на основе большой языковой модели лучше интерпретирует инструкции, но сохраняет высокий уровень нарушений из-за ошибок связывания текстового ограничения с состояниями и действиями среды.

Полученные результаты показывают, что безопасное обучение с подкреплением нельзя свести только к добавлению штрафа в целевую функцию. В задачах с отложенными последствиями требуется оценивать будущий накопленный штраф действий; в задачах большого горизонта~--- направлять исследование через промежуточные подцели; в задачах с текстовыми ограничениями~--- связывать правило безопасности с конкретными объектами, действиями и состояниями среды.

При этом остаются направления, требующие дальнейшего исследования. Для методов на основе моделей мира важна устойчивость при увеличении горизонта планирования и при переносе из симуляторов на реальные робототехнические системы. Для иерархических методов остаётся открытым вопрос перепланирования во время исполнения, особенно если условия среды меняются или выбранная подцель становится недостижимой. Для задач с естественно-языковыми ограничениями ключевой проблемой остаётся надёжное связывание языка и наблюдений: агент должен не только понимать текст правила, но и применять его к текущей ситуации в среде.

В целом диссертационная работа развивает безопасное обучение с подкреплением в трёх направлениях: оценка долгосрочного риска с помощью моделей мира, безопасное иерархическое исследование и постановка задач с естественно-языковыми целями и ограничениями. Эти результаты могут быть использованы при разработке автономных агентов, которым необходимо действовать в новых условиях, исследовать среду и при этом соблюдать заданные ограничения безопасности.
